% ============================================================
% thesis-main.tex — 北京大学博士学位论文主文件
% 基于 pkuthss v1.9.4 (2024/03/07) + Word模板精确校正
%
% 编译命令（完整流程）：
%   xelatex thesis-main
%   biber thesis-main
%   xelatex thesis-main
%   xelatex thesis-main
% 或：latexmk -xelatex thesis-main
% ============================================================

% ugly 选项：章标题三号(16pt)，小节标题四号(14pt) — 匹配Word模板
% nopkufont：禁用 unicode-math，避免 amssymb/algorithm2e 冲突
\documentclass[UTF8, ugly, nopkufont]{pkuthss}

% ============================================================
% 字体设置（Word模板：正文宋体+Times New Roman；标题黑体）
% ============================================================
\usepackage{fontspec}
\setmainfont{Times New Roman}        % 英文/数字：Times New Roman
\setsansfont{Arial}                  % 英文无衬线：Arial（ABSTRACT标题用）

% ============================================================
% 页面尺寸精确校正（对照Word模板实测值）
% pkuthss ugly 默认：top=3.1cm（应为3.0cm），headsep=0.6cm（应为0.5cm）
% 关系：header位置 = top - headheight - headsep = 3.0 - 0.5 - 0.5 = 2.0cm ✓
% ============================================================
\geometry{top=3.0cm, bottom=2.5cm, headsep=0.5cm, footskip=0.75cm}

% ============================================================
% 参考文献（GB/T 7714-2015 顺序编码制）
% ============================================================
\usepackage[
  backend = biber,
  style   = gb7714-2015,
  sorting = none
]{biblatex}
\renewcommand*{\bibfont}{\zihao{5}\linespread{1.27}\selectfont}
\setlength{\bibitemsep}{3bp}
\addbibresource{references/references.bib}

% ============================================================
% 常用宏包
% ============================================================
\usepackage{amsmath, amssymb, amsthm}
\usepackage[ruled, linesnumbered]{algorithm2e}
\usepackage{booktabs}                % 三线表
\usepackage{longtable}               % 跨页表格

% ============================================================
% 标题格式校正（精确匹配Word模板间距）
%
% Word模板实测（twips→pt）：
%   章：before=480/20=24pt, after=360/20=18pt, 单倍行距
%   一级节(heading 2)：before=480/20=24pt, after=120/20=6pt
%     （注：pkuthss默认before=20pt，Word实测24pt）
%   二级节(heading 3)：before=240/20=12pt, after=120/20=6pt
%   三级节(heading 4)：before=240/20=12pt, after=120/20=6pt
% ============================================================
\ctexset{
  chapter = {
    beforeskip = 24bp,
    afterskip  = 18bp plus 0.2ex,
    format     = {\linespread{1.0}\zihao{3}\bfseries\centering},
  },
  section = {
    beforeskip = 24bp,
    afterskip  = 6bp plus 0.2ex,
    format     = {\linespread{1.0}\zihao{4}\bfseries},
  },
  subsection = {
    beforeskip = 12bp plus 1ex minus 0.2ex,
    afterskip  = 6bp plus 0.2ex,
    format     = {\linespread{1.0}\fontsize{13bp}{15.6bp}\selectfont\bfseries},
  },
  subsubsection = {
    beforeskip = 12bp plus 1ex minus 0.2ex,
    afterskip  = 6bp plus 0.2ex,
    format     = {\linespread{1.0}\zihao{-4}\bfseries},
  },
}

% ============================================================
% 目录格式（章条目：黑体小四；节条目：宋体小四，段前0）
% ============================================================
\renewcommand{\cftchapfont}{\heiti}
\renewcommand{\cftchappagefont}{\heiti}
\setlength{\cftbeforesecskip}{0pt}
\setlength{\cftbeforesubsecskip}{0pt}

% ============================================================
% 英文摘要重定义（Word模板：仅"ABSTRACT"为章标题，不含英文正文题目）
% pkuthss原始版本在英文摘要页显示英文论文题目+作者信息，不符合北大Word模板
% ============================================================
\makeatletter
\renewenvironment{eabstract}{%
  \thss@int@pdfmark{\eabstractname}{eabstract}%
  \chapter*{\sffamily\eabstractname}\markboth{\eabstractname}{}%
}{%
  \vfill\noindent\textbf{\label@ekeywords}{\@ekeywords}%
}
\makeatother

% 双盲评审开关（盲审版改为 \blindtrue）
\newif\ifblind\blindfalse

% ============================================================
% 论文基本信息（TODO：填写实际信息）
% ============================================================
\pkuthssinfo{
  cthesisname  = {博士学位论文},
  ethesisname  = {Doctor Thesis},
  thesiscover  = {博士研究生学位论文},
  % 论文标题（中英文）
  ctitle       = {论文题目（不超过20字）},
  etitle       = {Thesis Title in English},
  % 作者信息
  cauthor      = {作者姓名},
  eauthor      = {Author Name},
  studentid    = {0000000000},
  % 院系与专业
  school       = {工学院},
  cmajor       = {生物医学工程},
  emajor       = {Biomedical Engineering},
  direction    = {生物信息学算法开发},
  % 导师（多位导师用 \\ 换行，mentorlines 相应增加）
  mentorlines  = {1},
  cmentor      = {导师姓名~~教授},
  ementor      = {Prof. Supervisor Name},
  % 关键词
  ckeywords    = {关键词1，关键词2，关键词3，关键词4，关键词5},
  ekeywords    = {Keyword One, Keyword Two, Keyword Three, Keyword Four, Keyword Five},
  % 日期（三号宋体汉字，如"二〇二五年六月"）
  date         = {二〇二五年六月},
  % 学位类型：1=学术学位，2=专业学位
  degreetype   = {1},
  % 双盲评审信息（非盲审可忽略）
  blindid      = {0000000000},
  discipline   = {生物医学工程},
}

% ============================================================
% 文档开始
% ============================================================
\begin{document}

% ------ 前置部分（封面、版权、摘要、目录）------
\frontmatter
\pagestyle{empty}

% 中文封面（自动生成）
\ifblind\makeblind\else\maketitle\fi

% 版权声明（正式提交时从研究生院门户打印替换本页，单面打印）
\cleardoublepage
\begin{center}
  \vspace*{8cm}
  {\zihao{3}\heiti 版权声明}\\[1.5cm]
  {\zihao{-4}\songti
    本论文版权归北京大学所有，使用须遵守学校相关规定。\\[0.5cm]
    （正式提交时请从研究生院门户下载打印版权声明页替换本页）
  }
\end{center}

% 摘要与目录（带页眉页脚，罗马数字页码）
\cleardoublepage
\pagestyle{plain}
\pagenumbering{Roman}

% 中文摘要（≥800字，关键词3~5个置底）
\begin{cabstract}
  % TODO：填写中文摘要（800-1000字）
  某某问题是生物信息学领域的重要研究方向。本文针对该问题，提出了……

  本文采用了……方法，主要工作包括：（1）……；（2）……；（3）……

  实验结果表明，本文方法在……数据集上取得了……的性能，较已有方法提升了……

  本研究的主要贡献如下：……
\end{cabstract}

% 英文摘要（与中文摘要内容一致，ABSTRACT居中三号）
\begin{eabstract}
  % TODO：填写英文摘要（内容与中文摘要一致）
  This thesis focuses on ...

  The main contributions of this work are as follows:
  (1) ...; (2) ...; (3) ...

  Experimental results demonstrate that the proposed method achieves
  state-of-the-art performance on ... datasets.
\end{eabstract}

% 目录（章标题居左黑体小四，节标题缩进宋体小四）
\tableofcontents

% ============================================================
% 正文
% ============================================================
\mainmatter

% ============================================================
% 第一章：引言 / 绪论
% ============================================================
\chapter{引言}
\label{ch:introduction}

% TODO: 本章写作目标
% - 研究背景与意义（500~800字）
% - 国内外研究现状综述（可合并至第二章）
% - 本文研究内容与主要贡献
% - 论文组织结构

\section{研究背景与意义}
\label{sec:background}

\section{国内外研究现状}
\label{sec:related-work}

\section{主要研究内容与贡献}
\label{sec:contributions}

\section{论文组织结构}
\label{sec:organization}

% ============================================================
% 第二章：文献综述
% ============================================================
\chapter{相关工作}
\label{ch:literature}

% TODO: 本章写作目标
% - 生物信息学算法发展脉络
% - 核心方法论综述（序列比对/变异检测/机器学习等）
% - 现有方法的局限性 → 引出本文动机

\section{生物信息学算法概述}
\label{sec:bioinformatics-overview}

\section{相关算法方法综述}
\label{sec:algorithm-review}

\section{现有方法的不足}
\label{sec:limitations}

\section{本章小结}
\label{sec:ch2-summary}

% ============================================================
% 第三章：方法（算法设计）
% ============================================================
\chapter{算法设计与理论分析}
\label{ch:methods}

% TODO: 本章写作目标
% - 问题形式化定义（数学符号表）
% - 核心算法设计（伪代码 + 流程图）
% - 时间/空间复杂度分析
% - 理论正确性证明（如适用）

\section{问题定义与形式化描述}
\label{sec:problem-definition}

\section{算法总体框架}
\label{sec:algorithm-overview}

\section{核心模块设计}
\label{sec:core-modules}

\section{复杂度分析}
\label{sec:complexity}

\section{本章小结}
\label{sec:ch3-summary}

% ============================================================
% 第四章：实验与结果分析
% ============================================================
\chapter{实验与结果分析}
\label{ch:experiments}

% TODO: 本章写作目标
% - 实验数据集描述（真实/模拟数据，来源，统计信息）
% - 基准方法（baselines）设置
% - 评估指标定义
% - 实验结果与对比分析
% - 消融实验（ablation study）

\section{实验设置}
\label{sec:exp-setup}

\subsection{数据集描述}
\label{subsec:datasets}

\subsection{基准方法与评估指标}
\label{subsec:baselines}

\subsection{实验环境}
\label{subsec:env}

\section{主实验结果}
\label{sec:main-results}

\section{消融实验}
\label{sec:ablation}

\section{参数敏感性分析}
\label{sec:sensitivity}

\section{本章小结}
\label{sec:ch4-summary}

% ============================================================
% 第五章：讨论
% ============================================================
\chapter{讨论}
\label{ch:discussion}

\section{结果解读与生物学意义}
\label{sec:interpretation}

\section{方法局限性分析}
\label{sec:limitations}

\section{与现有工作的对比}
\label{sec:comparison}

\section{本章小结}
\label{sec:ch5-summary}

% ============================================================
% 第六章：结论与展望
% ============================================================
\chapter{结论与展望}
\label{ch:conclusion}

\section{工作总结}
\label{sec:summary}

\section{主要创新点}
\label{sec:innovations}

\section{未来工作展望}
\label{sec:future-work}


% ============================================================
% 后置部分
% ============================================================

% 参考文献（五号宋体，行距固定16磅，段前3磅）
\printbibliography[heading = bibintoc, title = {参考文献}]

% 附录（有则取消注释）
% \appendix
% \include{appendices/appendix-a}

\backmatter

% 个人简历与学术成果
\specialchap{个人简历、在学期间发表的学术论文与研究成果}

\subsection*{个人简历}
% TODO：填写个人简历

\subsection*{科研成果}
% TODO：填写发表论文和研究成果（格式同参考文献）

% 致谢
\ifblind\else
\specialchap{致\quad 谢}
% TODO：填写致谢内容（不超过1000字）
\fi

% 原创性声明（正式提交时从门户打印后附于此处）
\specialchap{北京大学学位论文原创性声明和使用授权说明}
{\zihao{-4}\songti
  本人郑重声明：所呈交的学位论文，是本人在导师的指导下，独立进行研究工作所取得的成果。
  除文中已经注明引用的内容外，本论文不含任何其他个人或集体已经发表或撰写过的作品成果。
  对本文的研究做出重要贡献的个人和集体，均已在文中以明确方式标明。

  \vspace{2cm}
  论文作者签名：\hspace{4cm}日期：\hspace{1cm}年\hspace{0.8cm}月\hspace{0.8cm}日
}

\end{document}
